\scanpage{147}
induit donc une équivalence sur les $\Pi_1$. Mais le bât blesse, du fait qu'on n'a pas de canonicité --- sans compter la nécessité d'hypothèses \inserted{restrictives} gênantes. \deleted{\unsure{Pour trouver quelque chose de canonique}}
\diagram{\begin{tikzcd}[ampersand replacement=\&]
\deleted{\Pi_1^{\unsure{\mathrm{arcs}|\mathrm{accès}}}P_{D'^{*}}}\arrow[r]\arrow[d,"\deleted{\sim}"'] \& \deleted{\Pi_1^{\unsure{\mathrm{rev}}}\mathcal V_{D',Y}}\\
\deleted{\Pi_1^{\unsure{\mathrm{rev}}}P_{D'^{*}}} \& \deleted{\uline{NB}\;\Pi_1^{\unsure{\mathrm{arcs}|\mathrm{accès}}}\mathcal V}
\end{tikzcd}}
\margin{Mais on peut dire qu'on a $\mathcal V^{*}_{D'^{*}}\lhook\joinrel\longrightarrow\mathcal V^{*}_{D'}$ qui induit une équivalence sur les $\Pi_1$, donc il suffit de comparer les $\Pi_1$ des $\mathcal V^{*}_{D'^{*}}$ et $P_{D'^{*}}$, \unsure{en oubliant les généralités précédentes}.}

On a oublié de préciser d'ailleurs, en relation avec le diviseur $\Theta$, qu'on veut que l'é\deleted{\unsure{isom.}}\inserted{plongement (240)} soit tel que l'image inverse de $P^{\deleted{\unsure{*}}}_{D'^{*}}$, identifié à un ouvert de \inserted{($P_{D'}$, lui $=$ ouvert de)} $\check N_{D',Y}$, soit juste $\mathcal V^{*}_{D',Y}$, de sorte qu'on \unsure{trouve} \inserted{bien} \unsure{ainsi|aussi}
\[
\deleted{\Pi_1\mathcal V^{*}_{D',Y}\longrightarrow\Pi_1P_{D'^{*}}}
\]
\[
(241)\qquad\mathcal V^{*}_{D',Y}\longrightarrow P_{D'^{*}}
\]
d'où
\[
(242)\qquad\Pi_1\mathcal V^{*}_{D',Y}\longrightarrow\Pi_1P_{D'^{*}}.
\]
Pour avoir un morphisme de groupoïdes défini de façon plus ``canonique'', et
\margin{146}
